% Part: proof-theory
% Chapter: cut-elimination
% Section: interpolation

\documentclass[../../../include/open-logic-section]{subfiles}

\begin{document}

\olfileid{pt}{cut}{itp}

\olsection{لم مائهارا و قضیهٔ درونیابی کریگ}

قضیهٔ درونیابی، از ویلیام کریگ، می‌گوید اگر $!A\Entails !B$، آنگاه
!!a{sentence}~$!C$ای، موسوم به \emph{درونیاب}، وجود دارد که
$!A\Entails !C$ و $!C\Entails !B$. نکتهٔ مهم این است که می‌توان~$!C$
را چنان برگزید که همهٔ !!{constant}ها و !!{predicate}های ظاهرشده در
آن هم در~$!A$ و هم در~$!B$ ظاهر شوند. (فرض می‌کنیم هیچ !!{function}ی
در کار نیست.) بدون این محدودیت، خود~$!A$ و~$!B$ به‌سادگی می‌توانستند
درونیاب باشند.

قضیهٔ درونیابی برای منطق کلاسیک مرتبهٔ اول برقرار است و آن را
«آخرین ویژگی مهم منطق مرتبهٔ اول» که کشف شد نامیده‌اند. این قضیه در
نظریهٔ مدل و استدلال خودکار کاربردهای مهمی دارد. انگیزه و کاربرد اصلی
آن در فلسفهٔ علم بود، زیرا می‌توان با آن مفاهیم اولیه‌ای را بررسی کرد
که یک نظریه بر پایهٔ آن‌ها اصل‌موضوعی می‌شود یا نمی‌شود. این قضیه
همچنین قضیهٔ تعریف‌پذیری بث را نتیجه می‌دهد: اگر نظریه‌ای بتواند
مفهومی را به‌طور ضمنی تعریف کند، می‌تواند آن را به‌طور صریح نیز تعریف
کند.

مانند بسیاری از نتایج منطق ریاضی، قضیهٔ درونیابی را هم با روش‌های
نظریهٔ مدل و هم با روش‌های نظریهٔ اثبات می‌توان اثبات کرد. مزیت
روش‌های نظریه‌اثباتی این است که نه‌فقط وجود درونیاب را نشان می‌دهند،
بلکه الگوریتمی برای محاسبهٔ آن از !!a{proof}ی برای
$!A\Entails !B$، مثلاً !!a{proof}ی در~\Log{G3c} برای
$!A\Sequent !B$، فراهم می‌کنند.

بگذارید $\Lang L(!A)$ مجموعهٔ !!{constant}ها و !!{predicate}های
درون~$!A$ باشد و
$\Lang L(\Gamma)=\bigcup\Setabs{\Lang L(!A)}{!A\in\Gamma}$.
در سراسر این بخش فرض می‌کنیم با !!{formula}های بدون !!{function}
سروکار داریم.

\begin{thm}[قضیهٔ درونیابی کریگ]\ollabel{thm:interpolation}
بگذارید !!{language}~$!A$ برابر $\Lang L(!A)$ و زبان~$!B$ برابر
$\Lang L(!B)$ باشد. اگر $\Log{G3c}\Proves !A\Sequent !B$، آنگاه
!!a{formula}~$!C$ای در !!{language}
$\Lang L(!C)\subseteq\Lang L(!A)\cap\Lang L(!B)$ وجود دارد که
$\Log{G3c}\Proves !A\Sequent !C$ و
$\Log{G3c}\Proves !C\Sequent !B$.
\end{thm}

برای بهره‌گرفتن از ساختار !!{proof}ها در~\Log{G3c}، تعمیمی از قضیهٔ
درونیابی را اثبات می‌کنیم. مقدم و تالی دنباله‌ها را هرکدام به دو بخش
تقسیم‌شده تصور کنید. می‌نویسیم
$\maeh{\Gamma_1}{\Gamma_2}\Sequent
\maeh{\Delta_1}{\Delta_2}$ تا نشان دهیم $\Gamma_1$ و~$\Delta_1$ به
بخش نخست و $\Gamma_2$ و~$\Delta_2$ به بخش دوم تعلق دارند. قضیهٔ
درونیابی بی‌درنگ از لم زیر نتیجه می‌شود.

\begin{lem}[لم مائهارا]\ollabel{lem:maehara}
اگر
$\Log{G3c}\Proves\maeh{\Gamma_1}{\Gamma_2}\Sequent
\maeh{\Delta_1}{\Delta_2}$، آنگاه !!a{formula}~$!C$ای در
!!{language}
\[
\Lang L(!C)\subseteq
\Lang L(\Gamma_1,\Delta_1)\cap\Lang L(\Gamma_2,\Delta_2)
\]
وجود دارد که
$\Log{G3c}\Proves\Gamma_1\Sequent\Delta_1,!C$ و
$\Log{G3c}\Proves !C,\Gamma_2\Sequent\Delta_2$.
\end{lem}

\begin{proof}
فرض کنید
$\pi\Proves\maeh{\Gamma_1}{\Gamma_2}\Sequent
\maeh{\Delta_1}{\Delta_2}$. ادعا را با استقرا بر
$n=\pheight{\pi}$ اثبات می‌کنیم.

اگر $n=0$، آنگاه
$\maeh{\Gamma_1}{\Gamma_2}\Sequent
\maeh{\Delta_1}{\Delta_2}$ یک اصل است و !!{formula} اتمی~$!A$ای هم
در $\Gamma_1,\Gamma_2$ و هم در $\Delta_1,\Delta_2$ ظاهر می‌شود. برحسب
اینکه $!A$ در سمت چپ و راست به بخش نخست یا دوم تعلق دارد چهار حالت
هست:
\begin{enumerate}
\item $!A\in\Gamma_1$ و $!A\in\Delta_1$. می‌خواهیم~$!C$ای بیابیم که
هر دو $\Gamma_1\Sequent\Delta_1,!C$ و
$!C,\Gamma_2\Sequent\Delta_2$ !!{provable} باشند. دنبالهٔ نخست،
صرف‌نظر از انتخاب~$!C$، از پیش اصل است. تنها انتخابی که تضمین می‌کند
دنبالهٔ دوم !!{provable} باشد، $!C\ident\lfalse$ است.
\item $!A\in\Gamma_1$ و $!A\in\Delta_2$. هر دو دنبالهٔ
$\Gamma_1\Sequent\Delta_1,!A$ و
$!A,\Gamma_2\Sequent\Delta_2$ اصل‌اند؛ پس $!C\ident !A$ را
برمی‌گزینیم.
\item $!A\in\Gamma_2$ و $!A\in\Delta_1$. دنباله‌های
$!A,\Gamma_1\Sequent\Delta_1$ و
$\Gamma_2\Sequent\Delta_2,!A$ اصل‌اند، اما~$!A$ در سمت نادرست قرار
دارد و درونیاب نیست. با استفاده از \RightR{\lnot} و
\LeftR{\lnot} می‌توانیم
$\Gamma_1\Sequent\Delta_1,\lnot !A$ و
$\lnot !A,\Gamma_2\Sequent\Delta_2$ را !!{prove} کنیم.
\item $!A\in\Gamma_2$ و $!A\in\Delta_2$. تمرین.
\end{enumerate}
حالت‌هایی که دنباله به سبب
$\lfalse\in\Gamma_1$ یا $\lfalse\in\Gamma_2$ اصل است به‌عنوان تمرین
واگذار می‌شوند.

اگر $n>0$، $\pi$ به یک استنتاج منطقی پایان می‌یابد. باید حالت‌ها را
برحسب قاعدهٔ به‌کاررفته و بخشی که فرمول اصلی به آن تعلق دارد جدا کنیم.
در هر حالت از فرض استقرا استفاده می‌کنیم: لم برای هر !!{proof}
~$\pi_1$ با $\pheight{\pi_1}<n$ و برای هر تقسیم دنبالهٔ پایانی آن به
دو بخش برقرار است. چند حالت را بررسی کنیم.
\begin{enumerate}
\item $\pi$ به \RightR{\lnot} با فرمول اصلی~$\lnot !A$ پایان می‌یابد.
برحسب اینکه $\lnot !A$ به بخش نخست یا دوم تعلق دارد، $\pi$ به یکی از
دو شکل زیر پایان می‌یابد:
\[
\AxiomC{}
\RightLabel{$\pi_1$}
\Deduce$\maeh{!A, \Gamma_1}{\Gamma_2} \fCenter
  \maeh{\Delta_1}{\Delta_2}$
\RightLabel{\RightR{\lnot}}
\UnaryInf$\maeh{\Gamma_1}{\Gamma_2} \fCenter
  \maeh{\Delta_1, \lnot !A}{\Delta_2}$
\DisplayProof
\]
یا
\[
\AxiomC{}
\RightLabel{$\pi_1$}
\Deduce$\maeh{\Gamma_1}{!A, \Gamma_2} \fCenter
  \maeh{\Delta_1}{\Delta_2}$
\RightLabel{\RightR{\lnot}}
\UnaryInf$\maeh{\Gamma_1}{\Gamma_2} \fCenter
  \maeh{\Delta_1}{\Delta_2, \lnot !A}$
\DisplayProof
\]
اگر !!{formula} اصلی~$\lnot !A$ به بخش نخست تعلق دارد، !!{formula}
فعال~$!A$ را در مقدمه نیز به بخش نخست داده‌ایم. این انتخاب در اختیار
ماست. فرض استقرا با قراردادن فرمول فعال در بخش دیگر نیز اعمال می‌شد،
اما آنگاه درونیابی نمی‌یافتیم.

نخست وضعیتی را در نظر بگیرید که $\lnot !A$ به بخش نخست تعلق دارد.
بنا بر فرض استقرا، دنبالهٔ پایانی~$\pi_1$ درونیابی، یعنی
!!a{formula}~$!C_1$، دارد که:
\begin{align*}
!A, \Gamma_1 & \Sequent \Delta_1, !C_1 &&\text{و} &
!C_1, \Gamma_2 & \Sequent \Delta_2
\intertext{اثبات‌پذیرند. به !!a{formula}~$!C$ای نیاز داریم که هر دو
دنبالهٔ زیر !!{provable} باشند:}
\Gamma_1 & \Sequent \Delta_1, \lnot !A, !C &&\text{و} &
!C, \Gamma_2 & \Sequent \Delta_2
\end{align*}
کافی است $!C\ident !C_1$ بگیریم: فرض استقرا گفته است دنبالهٔ راست
!!{provable} است و دنبالهٔ چپ با~$\RightR{\lnot}$ از دنبالهٔ فرض استقرا
به دست می‌آید. همچنین
$\Lang L(!C_1)\subseteq
\Lang L(!A,\Gamma_1,\Delta_1)\cap\Lang L(\Gamma_2,\Delta_2)$.
چون $\Lang L(\lnot !A)=\Lang L(!A)$ و
$\Lang L(!C)=\Lang L(!C_1)$، داریم
\[
\Lang L(!C_1)\subseteq
\Lang L(\Gamma_1,\Delta_1,\lnot !A)\cap
\Lang L(\Gamma_2,\Delta_2).
\]

در وضعیت دوم $\lnot !A$ به بخش دوم تعلق دارد. !!{formula} فعال در
مقدمه را نیز به بخش دوم می‌دهیم و فرض استقرا را اعمال می‌کنیم:
\begin{align*}
\Gamma_1 & \Sequent \Delta_1, !C_1 &&\text{و} &
!C_1, !A, \Gamma_2 & \Sequent \Delta_2
\intertext{اثبات‌پذیرند. با اعمال \RightR{\lnot} بر دومی،
!!{proof}هایی برای این دو دنباله به دست می‌آوریم:}
\Gamma_1 & \Sequent \Delta_1, !C_1 &&\text{و} &
!C_1, \Gamma_2 & \Sequent \Delta_2, \lnot !A
\end{align*}
این‌ها دقیقاً همان دنباله‌هایی هستند که باید !!{provable} باشند؛ پس باز
$!C\ident !C_1$. مانند پیش،
$\Lang L(!C_1)\subseteq
\Lang L(\Gamma_1,\Delta_1)\cap
\Lang L(\Gamma_2,\Delta_2,\lnot !A)$.

\item $\pi$ به \RightR{\lif} با !!{formula} اصلی~$!A\lif !B$ پایان
می‌یابد. باز دو حالت داریم و !!{proof}~$\pi$ به یکی از دو شکل زیر پایان می‌یابد:
\[
\AxiomC{}
\RightLabel{$\pi_1$}
\Deduce$\maeh{!A, \Gamma_1}{\Gamma_2} \fCenter
  \maeh{\Delta_1, !B}{\Delta_2}$
\RightLabel{\RightR{\lif}}
\UnaryInf$\maeh{\Gamma_1}{\Gamma_2} \fCenter
  \maeh{\Delta_1, !A \lif !B}{\Delta_2}$
\DisplayProof
\]
یا
\[
\AxiomC{}
\RightLabel{$\pi_1$}
\Deduce$\maeh{\Gamma_1}{!A, \Gamma_2} \fCenter
  \maeh{\Delta_1}{\Delta_2, !B}$
\RightLabel{\RightR{\lif}}
\UnaryInf$\maeh{\Gamma_1}{\Gamma_2} \fCenter
  \maeh{\Delta_1}{\Delta_2, !A \lif !B}$
\DisplayProof
\]
باز !!{formula}های فعال را در مقدمه به همان بخشی داده‌ایم که
!!{formula} اصلی در نتیجه دارد. در وضعیت نخست، فرض استقرا
!!{proof}هایی برای دنباله‌های زیر می‌دهد:
\begin{align*}
!A, \Gamma_1 & \Sequent \Delta_1, !B, !C_1 &&\text{و} &
!C_1, \Gamma_2 & \Sequent \Delta_2
\intertext{دنبالهٔ چپ مقدمهٔ مناسبی برای \RightR{\lif} است؛ پس
دنباله‌های !!{provable} زیر را داریم:}
\Gamma_1 & \Sequent \Delta_1, !A \lif !B, !C_1 &&\text{و} &
!C_1, \Gamma_2 & \Sequent \Delta_2
\end{align*}
پس $!C_1$ برای دنبالهٔ پایانی~$\pi$ نیز درونیاب است و دوباره
$!C\ident !C_1$. وضعیت دیگر نیز به همین روش حل می‌شود.

\item $\pi$ به \LeftR{\lif} با !!{formula} اصلی~$!A\lif !B$ پایان
می‌یابد. اگر $!A\lif !B$ به بخش نخست تعلق داشته باشد، !!{proof}~$\pi$ چنین
پایان می‌یابد:
\[
\AxiomC{}
\RightLabel{$\pi_1$}
\Deduce$\maeh{\Gamma_1}{\Gamma_2} \fCenter
  \maeh{\Delta_1, !A}{\Delta_2}$
\AxiomC{}
\RightLabel{$\pi_2$}
\Deduce$\maeh{!B, \Gamma_1}{\Gamma_2} \fCenter
  \maeh{\Delta_1}{\Delta_2}$
\RightLabel{\LeftR{\lif}}
\BinaryInf$\maeh{!A \lif !B, \Gamma_1}{\Gamma_2} \fCenter
  \maeh{\Delta_1}{\Delta_2}$
\DisplayProof
\]
فرض استقرا چهار دنبالهٔ !!{provable} می‌دهد؛ دو دنباله برای
درونیاب~$!C_1$ مقدمهٔ چپ و دو دنباله برای درونیاب~$!C_2$ مقدمهٔ دوم:
\begin{align*}
\Gamma_1 & \Sequent \Delta_1, !A, !C_1 &&\text{(۱)} &
!C_1, \Gamma_2 & \Sequent \Delta_2 &&\text{(۲)}\\
!B, \Gamma_1 & \Sequent \Delta_1, !C_2 &&\text{(۳)} &
!C_2, \Gamma_2 & \Sequent \Delta_2 &&\text{(۴)}
\end{align*}
دنباله‌های (۱) و~(۳)، پس از تضعیف و افزودن~$!C_2$ به راست (۱) و
~$!C_1$ به راست (۳)، مقدمات \LeftR{\lif} می‌شوند:
\[
\AxiomC{}
\Deduce$\Gamma_1 \fCenter \Delta_1, !A, !C_1, !C_2$
\AxiomC{}
\Deduce$!B, \Gamma_1 \fCenter \Delta_1, !C_1, !C_2$
\RightLabel{\LeftR{\lif}}
\BinaryInf$!A \lif !B, \Gamma_1 \fCenter \Delta_1, !C_1, !C_2$
\DisplayProof
\]
با اعمال $\RightR{\lor}$ دنبالهٔ
\[
!A \lif !B, \Gamma_1 \fCenter \Delta_1, !C_1 \lor !C_2
\]
را به دست می‌آوریم. پس $!C_1\lor !C_2$ نامزد درونیاب است. دنبالهٔ
دیگر لازم را نیز از (۲) و~(۴) !!{prove} می‌کنیم:
\[
\AxiomC{}
\Deduce$!C_1, \Gamma_2 \fCenter \Delta_2$
\AxiomC{}
\Deduce$!C_2, \Gamma_2 \fCenter \Delta_2$
\RightLabel{\LeftR{\lor}}
\BinaryInf$!C_1 \lor !C_2, \Gamma_2 \fCenter \Delta_2$
\DisplayProof
\]
فقط باید زبان مناسب را بررسی کنیم. فرض استقرا می‌دهد:
\begin{align*}
\Lang L(!C_1) & \subseteq
\Lang L(\Gamma_1, \Delta_1, !A) \cap
\Lang L(\Gamma_2, \Delta_2) &\text{و}\\
\Lang L(!C_2) & \subseteq
\Lang L(!B, \Gamma_1, \Delta_1) \cap
\Lang L(\Gamma_2, \Delta_2)
\intertext{از آنجا که}
\Lang L(!A) & \subseteq \Lang L(!A \lif !B) &\text{و}\\
\Lang L(!B) & \subseteq \Lang L(!A \lif !B)
\intertext{هر دو رابطهٔ زیر برقرارند:}
\Lang L(!C_1) & \subseteq
\Lang L(\Gamma_1, \Delta_1, !A \lif !B) \cap
\Lang L(\Gamma_2, \Delta_2) &\text{و}\\
\Lang L(!C_2) & \subseteq
\Lang L(\Gamma_1, \Delta_1, !A \lif !B) \cap
\Lang L(\Gamma_2, \Delta_2)
\intertext{و در نتیجه، چون
$\Lang L(!C_1\lor !C_2)={}$}
\Lang L(!C_1)\cup\Lang L(!C_2) & \subseteq
\Lang L(\Gamma_1,\Delta_1,!A\lif !B)\cap
\Lang L(\Gamma_2,\Delta_2).
\end{align*}
این همان رابطهٔ مطلوب است.

اگر $!A\lif !B$ به بخش دوم تعلق داشته باشد، وضعیت متفاوت اما روش
مشابه است. فرض استقرا باز چهار دنبالهٔ !!{provable} می‌دهد:
\begin{align*}
\Gamma_1 & \Sequent \Delta_1, !C_1 &&\text{(۱)} &
!C_1, \Gamma_2 & \Sequent \Delta_2, !A &&\text{(۲)}\\
\Gamma_1 & \Sequent \Delta_1, !C_2 &&\text{(۳)} &
!C_2, !B, \Gamma_2 & \Sequent \Delta_2 &&\text{(۴)}
\end{align*}
اکنون دنباله‌های (۲) و~(۴)، همراه با چند !!{formula} افزوده‌شده با
تضعیف، مقدمات \LeftR{\lif} می‌شوند:
\[
\AxiomC{}
\Deduce$!C_1, !C_2, \Gamma_2 \fCenter \Delta_2, !A$
\AxiomC{}
\Deduce$!B, !C_1, !C_2, \Gamma_2 \fCenter \Delta_2$
\RightLabel{\LeftR{\lif}}
\BinaryInf$!A \lif !B, !C_1, !C_2, \Gamma_2 \fCenter \Delta_2$
\DisplayProof
\]
این بار یک !!{formula} ساخته‌شده از~$!C_1$ و~$!C_2$ را در مقدم لازم
داریم. با $\LeftR{\land}$ این کار انجام می‌شود و باید
$!C\ident(!C_1\land !C_2)$ بگیریم. خوشبختانه از دنباله‌های باقی‌ماندهٔ
(۱) و~(۳) نیز دنبالهٔ متناظر بخش دیگر را !!{prove} می‌کنیم:
\[
\AxiomC{}
\Deduce$\Gamma_1 \fCenter \Delta_1, !C_1$
\AxiomC{}
\Deduce$\Gamma_1 \fCenter \Delta_1, !C_2$
\RightLabel{\RightR{\land}}
\BinaryInf$\Gamma_1 \fCenter \Delta_1, !C_1 \land !C_2$
\DisplayProof
\]
مانند پیش،
$\Lang L(!C_1\land !C_2)\subseteq
\Lang L(\Gamma_1,\Delta_1)\cap
\Lang L(\Gamma_2,\Delta_2,!A\lif !B)$.

\item $\pi$ به \RightR{\lforall} با !!{formula} اصلی
~$\lforall[x][!A(x)]$ پایان می‌یابد:
\[
\AxiomC{}
\RightLabel{$\pi_1$}
\Deduce$\maeh{\Gamma_1}{\Gamma_2} \fCenter
  \maeh{\Delta_1, !A(c)}{\Delta_2}$
\RightLabel{\RightR{\lforall}}
\UnaryInf$\maeh{\Gamma_1}{\Gamma_2} \fCenter
  \maeh{\Delta_1, \lforall[x][!A(x)]}{\Delta_2}$
\DisplayProof
\]
یا
\[
\AxiomC{}
\RightLabel{$\pi_1$}
\Deduce$\maeh{\Gamma_1}{\Gamma_2} \fCenter
  \maeh{\Delta_1}{\Delta_2, !A(c)}$
\RightLabel{\RightR{\lforall}}
\UnaryInf$\maeh{\Gamma_1}{\Gamma_2} \fCenter
  \maeh{\Delta_1}{\Delta_2, \lforall[x][!A(x)]}$
\DisplayProof
\]
در وضعیت نخست، فرض استقرا !!{proof}هایی برای
$\Gamma_1\Sequent\Delta_1,!A(c),!C_1$ و
$!C_1,\Gamma_2\Sequent\Delta_2$ می‌دهد. با اعمال
$\RightR{\lforall}$ بر نخستین، به
$\Gamma_1\Sequent\Delta_1,\lforall[x][!A(x)],!C_1$ می‌رسیم.
شرط متغیر ویژه برقرار است، زیرا در قاعدهٔ \RightR{\lforall} پایانی
~$\pi$ برقرار بود. پس $!C_1$ درونیاب است، مشروط بر اینکه
\[
\Lang L(!C_1)\subseteq
\Lang L(\Gamma_1,\Delta_1,\lforall[x][!A(x)])\cap
\Lang L(\Gamma_2,\Delta_2).
\]
فرض استقرا رابطهٔ مشابه با~$!A(c)$ را می‌دهد. تنها نگرانی~$c$ است،
اما $c$ نمی‌تواند در~$!C_1$ ظاهر شود. اگر ظاهر می‌شد، هم در
$\Lang L(\Gamma_1,\Delta_1,!A(c))$ و هم در
$\Lang L(\Gamma_2,\Delta_2)$ بود؛ در نتیجه در نتیجهٔ استنتاج
\RightR{\lforall} در~$\pi$ ظاهر می‌شد و شرط متغیر ویژه را نقض می‌کرد.
وضعیت دیگر مشابه است.

\item $\pi$ به \RightR{\lexists} با !!{formula} اصلی
~$\lexists[x][!A(x)]$ پایان می‌یابد. باز دو وضعیت داریم. حالتی را
بررسی می‌کنیم که فرمول وجودی به بخش نخست تعلق دارد و دیگری را تمرین
می‌گذاریم. در وضعیت نخست، $\pi$ چنین پایان می‌یابد:
\[
\AxiomC{}
\RightLabel{$\pi_1$}
\Deduce$\maeh{\Gamma_1}{\Gamma_2} \fCenter
  \maeh{\Delta_1, \lexists[x][!A(x)], !A(t)}{\Delta_2}$
\RightLabel{\RightR{\lexists}}
\UnaryInf$\maeh{\Gamma_1}{\Gamma_2} \fCenter
  \maeh{\Delta_1, \lexists[x][!A(x)]}{\Delta_2}$
\DisplayProof
\]
فرض استقرا !!{proof}هایی برای
$\Gamma_1\Sequent\Delta_1,\lexists[x][!A(x)],!A(t),!C_1$ و
$!C_1,\Gamma_2\Sequent\Delta_2$ می‌دهد. با اعمال
$\RightR{\lexists}$ بر نخستین، به
$\Gamma_1\Sequent\Delta_1,\lexists[x][!A(x)],!C_1$ می‌رسیم. پس
$!C_1$ درونیاب است اگر:
\[
\Lang L(!C_1)\subseteq
\Lang L(\Gamma_1,\Delta_1,\lexists[x][!A(x)])\cap
\Lang L(\Gamma_2,\Delta_2).
\]
فرض استقرا همین رابطه را با~$!A(t)$ افزوده به زبان نخست می‌دهد. چون
!!{function} نداریم، ترم~$t$ یا متغیر است یا !!a{constant}. متغیر
عضو زبان نیست و مشکلی ایجاد نمی‌کند. اگر $t\ident b$ ثابت باشد نیز،
مگر در حالتی که $b\in\Lang L(!C_1)$ اما $b$ در اشتراک زبان‌های دو
بخش نباشد، $!C_1$ را بی‌تغییر می‌گیریم. فقط در همین حالتِ ثابت
نامشترک باید آن را انتزاع کنیم.

با جایگزینی همهٔ رخدادهای~$b$ در~$!C_1$ با !!a{variable}~$y$،
می‌نویسیم $!C_1\ident\Subst{!C_1'}{b}{y}$، که $!C_1'$ شامل~$b$ نیست.
افزون بر این، یا
$b\notin\Lang L(\Gamma_1,\Delta_1,\lexists[x][!A(x)])$ یا
$b\notin\Lang L(\Gamma_2,\Delta_2)$. به‌ترتیب می‌گیریم
$!C\ident\lforall[y][!C_1'(y)]$ یا
$!C\ident\lexists[y][!C_1'(y)]$. در حالت نخست:
\[
\AxiomC{}
\Deduce$\Gamma_1 \fCenter
  \Delta_1, \lexists[x][!A(x)], !A(b), !C_1'(b)$
\RightLabel{\RightR{\lexists}}
\UnaryInf$\Gamma_1 \fCenter
  \Delta_1, \lexists[x][!A(x)], !C_1'(b)$
\RightLabel{\RightR{\lforall}}
\UnaryInf$\Gamma_1 \fCenter
  \Delta_1, \lexists[x][!A(x)], \lforall[y][!C_1'(y)]$
\DisplayProof
\]
و
\[
\AxiomC{}
\Deduce$!C_1'(b), \lforall[y][!C_1'(y)], \Gamma_2 \fCenter \Delta_2$
\RightLabel{\LeftR{\lforall}}
\UnaryInf$\lforall[y][!C_1'(y)], \Gamma_2 \fCenter \Delta_2$
\DisplayProof
\]
!!{proof}های دنباله‌های بالایی از فرض استقرا می‌آیند؛ در دومی از
پذیرفتنی‌بودن تضعیف با $\lforall[y][!C_1'(y)]$ در مقدم نیز استفاده
می‌کنیم. شرط متغیر ویژهٔ \RightR{\lforall} در !!{proof} نخست برقرار
است، زیرا
$b\notin\Lang L(\Gamma_1,\Delta_1,\lexists[x][!A(x)])$ و~$!C_1'$
چنان تعریف شد که شامل~$b$ نباشد.

در حالت دوم، از فرض استقرا و پذیرفتنی‌بودن تضعیف با
~$\lexists[y][!C_1'(y)]$ داریم:
\[
\AxiomC{}
\Deduce$\Gamma_1 \fCenter
  \Delta_1, \lexists[x][!A(x)], !A(b),
  \lexists[y][!C_1'(y)], !C_1'(b)$
\RightLabel{\RightR{\lexists}}
\UnaryInf$\Gamma_1 \fCenter
  \Delta_1, \lexists[x][!A(x)],
  \lexists[y][!C_1'(y)], !C_1'(b)$
\RightLabel{\RightR{\lexists}}
\UnaryInf$\Gamma_1 \fCenter
  \Delta_1, \lexists[x][!A(x)], \lexists[y][!C_1'(y)]$
\DisplayProof
\]
و
\[
\AxiomC{}
\Deduce$!C_1'(b), \Gamma_2 \fCenter \Delta_2$
\RightLabel{\LeftR{\lexists}}
\UnaryInf$\lexists[y][!C_1'(y)], \Gamma_2 \fCenter \Delta_2$
\DisplayProof
\]
شرط متغیر ویژهٔ \LeftR{\lexists} در !!{proof} دوم برقرار است، زیرا
$b\notin\Lang L(\Gamma_2,\Delta_2)$ و~$!C_1'$ شامل~$b$ نیست.
\end{enumerate}
حالت‌های دیگر مشابه‌اند و آن‌ها را به‌عنوان تمرین واگذار می‌کنیم.
\end{proof}

\begin{prob}
فرض کنید رابط~$\lnand$ با قواعد زیر داشته باشیم:
\[
\Axiom$\Gamma \fCenter \Delta, !A$
\Axiom$\Gamma \fCenter \Delta, !B$
\RightLabel{\LeftR{\lnand}}
\BinaryInf$!A \lnand !B, \Gamma \fCenter \Delta$
\DisplayProof
\quad
\Axiom$!A, !B, \Gamma \fCenter \Delta$
\RightLabel{\RightR{\lnand}}
\UnaryInf$\Gamma \fCenter \Delta, !A \lnand !B$
\DisplayProof
\]
با انجام حالت‌های اثبات \olref{lem:maehara} که در آن‌ها $\pi$ به یکی
از این قواعد پایان می‌یابد، نشان دهید لم مائهارا همچنان برقرار است.
\end{prob}

اکنون که لم مائهارا ثابت شده است، از آن برای اثبات قضیهٔ درونیابی
کریگ استفاده می‌کنیم.

\begin{proof}[اثبات \olref{thm:interpolation}]
فرض کنید $!A\Sequent !B$ !!{provable} باشد. \olref{lem:maehara} را
بر $\maeh{!A}{\ }\Sequent\maeh{\ }{!B}$ اعمال کنید تا درونیاب~$!C$
به دست آید. آنگاه
$\Proves !A\Sequent !C$ و $\Proves !C\Sequent !B$.
\end{proof}

فرض کنید $\Gamma$ نظریه‌ای، یعنی مجموعه‌ای از !!{sentence}ها، در
!!a{language}~$\Lang L$ باشد که یک !!{predicate}~$n$موضعی~$R$ دارد.
می‌گوییم $\Gamma$، $R$ را \emph{به‌طور ضمنی تعریف می‌کند} اگر و تنها
اگر:
\begin{align*}
\Gamma, \Gamma'
& \Entails \lforall[x_1][\dots\lforall[x_n][
  (R(x_1,\dots,x_n)\liff R'(x_1,\dots,x_n))]],
\intertext{که در آن $\Gamma'$ از جایگزینی همهٔ رخدادهای~$R$ در
$\Gamma$ با~$R'\notin\Lang L$ حاصل می‌شود. $\Gamma$، $R$ را
\emph{به‌طور صریح تعریف می‌کند} اگر و تنها اگر
$!A(x_1,\dots,x_n)$ای بدون~$R$ وجود داشته باشد به‌طوری‌که}
\Gamma & \Entails
\lforall[x_1][\dots\lforall[x_n][
  (R(x_1,\dots,x_n)\liff !A(x_1,\dots,x_n))]].
\end{align*}
آشکار است اگر $\Gamma$، $R$ را صریحاً تعریف کند، آن را ضمنی نیز تعریف
می‌کند. عکس آن بدیهی نیست، اما برقرار است:

\begin{lem}[قضیهٔ تعریف‌پذیری بث]
اگر $\Gamma$، $R$ را به‌طور ضمنی تعریف کند، آن را به‌طور صریح نیز
تعریف می‌کند.
\end{lem}

\begin{proof}
فرض کنید $\Gamma$، $R$ را ضمنی تعریف می‌کند. $R'$ و~$\Gamma'$ را
مانند بالا بگیرید و
$\Lang L'=\Lang L(\Gamma')=(\Lang L\setminus\{R\})\cup\{R'\}$.
بنا بر فرض:
\begin{align*}
\Gamma,\Gamma' &
\Sequent \lforall[x_1][\dots\lforall[x_n][
  (R(x_1,\dots,x_n)\liff R'(x_1,\dots,x_n))]]
\intertext{بگذارید $c_1$، \dots،~$c_n$ !!{constant}هایی باشند که در
~$\Gamma$ ظاهر نمی‌شوند. بنا بر لم وارون‌پذیری:}
\Gamma,\Gamma',R(c_1,\dots,c_n) &
\Sequent R'(c_1,\dots,c_n)
\intertext{لم مائهارا را بر دنبالهٔ زیر اعمال کنید:}
\maeh{\Gamma,R(c_1,\dots,c_n)}{\Gamma'} &
\Sequent \maeh{\ }{R'(c_1,\dots,c_n)}
\intertext{تا $!A\ident !A(x_1,\dots,x_n)$ای بیابید که دنباله‌های زیر}
\Gamma,R(c_1,\dots,c_n) &
\Sequent !A(c_1,\dots,c_n)\\
!A(c_1,\dots,c_n),\Gamma' &
\Sequent R'(c_1,\dots,c_n)
\intertext{!!{proof} داشته باشند. لم مائهارا به‌درستی می‌دهد
$\Lang L(!A(c_1,\dots,c_n))\subseteq{}$}
\Lang L(\Gamma,R(c_1,\dots,c_n)) &\cap
\Lang L(\Gamma',R'(c_1,\dots,c_n))\\
&=
(\Lang L(\Gamma)\setminus\{R\})\cup\{c_1,\dots,c_n\}.
\intertext{پس پس از انتزاع ثابت‌های تازه به متغیرها،
$!A(x_1,\dots,x_n)$ شامل نه~$R$ و نه~$R'$ است. در !!{proof} دنبالهٔ
دوم، همه‌جا $R'$ را با~$R$ جایگزین کنید تا !!a{proof}ی برای دنبالهٔ
زیر به دست آید:}
!A(c_1,\dots,c_n),\Gamma &
\Sequent R(c_1,\dots,c_n).
\end{align*}
از دنبالهٔ نخست با \RightR{\lif}،
$\Gamma\Sequent R(\vec c)\lif !A(\vec c)$، و از دنبالهٔ دوم،
$\Gamma\Sequent !A(\vec c)\lif R(\vec c)$ را به دست می‌آوریم. دو جهت
را با قواعد گزاره‌ای در یک دوشرطی ترکیب می‌کنیم:
\[
\Gamma\Sequent
R(c_1,\dots,c_n)\liff !A(c_1,\dots,c_n).
\]
سپس \RightR{\lforall} را برای $c_1,\dots,c_n$ اعمال می‌کنیم و
!!a{proof}ی برای
\[
\Gamma \Sequent
\lforall[x_1][\dots\lforall[x_n][
  (R(x_1,\dots,x_n)\liff !A(x_1,\dots,x_n))]]
\]
به دست می‌آوریم. این نشان می‌دهد $!A(x_1,\dots,x_n)$، $R$ را در
~$\Gamma$ به‌طور صریح تعریف می‌کند.
\end{proof}

نتیجهٔ سومی که با درونیابی هم‌ارز است و مانند آن با لم مائهارا اثبات
می‌شود چنین است: اگر دو نظریهٔ سازگار~$\Gamma_1$ و~$\Gamma_2$ یک
نظریهٔ کامل~$\Gamma$ را در زبان
$\Lang L(\Gamma)\subseteq
\Lang L(\Gamma_1)\cap\Lang L(\Gamma_2)$ در بر داشته باشند، آنگاه
$\Gamma_1\cup\Gamma_2$ سازگار است.

به یاد آورید نظریه‌ای کامل است اگر برای هر !!{sentence}~$!A$ در
زبانش، یا $\Gamma\Proves !A$ یا $\Gamma\Proves\lnot !A$. چون
$\Gamma_1$ و~$\Gamma_2$ بسط‌های سازگار~$\Gamma$ هستند، خود~$\Gamma$
سازگار است. در نتیجه، اگر $!A$ !!a{sentence}‌ای در زبان مشترک
$\Lang L(\Gamma_1)\cap\Lang L(\Gamma_2)$ باشد، نمی‌توان هم
$\Gamma_1\Entails !A$ و هم
$\Gamma_2\Entails\lnot !A$ داشت. برای دیدن این امر، خلافش را فرض
کنید. از کامل‌بودن~$\Gamma$، یا $\Gamma\Entails !A$ یا
$\Gamma\Entails\lnot !A$. حالت دوم با
$\Gamma\subseteq\Gamma_1$ ناسازگاری~$\Gamma_1$ را نتیجه می‌دهد؛ پس
$\Gamma\Entails !A$. آنگاه $\Gamma\subseteq\Gamma_2$ می‌دهد
$\Gamma_2\Entails !A$، که همراه با
$\Gamma_2\Entails\lnot !A$ ناسازگاری~$\Gamma_2$ را نتیجه می‌دهد.

برای اثبات فقط کافی است هیچ !!a{sentence}~$!A$ای در
!!{language} مشترک
$\Lang L(\Gamma_1)\cap\Lang L(\Gamma_2)$ وجود نداشته باشد که
$\Gamma_1\Entails !A$ و $\Gamma_2\Entails\lnot !A$. اکنون اثبات
نظریه‌اثباتی آن را با لم مائهارا می‌آوریم.

\begin{thm}[قضیهٔ سازگاری مشترک رابینسون]
فرض کنید $\Gamma_1$ و~$\Gamma_2$ نظریه‌هایی سازگارند و هیچ~$!A$ای در
!!{language}
$\Lang L(\Gamma_1)\cap\Lang L(\Gamma_2)$ نیست که هم
$\Gamma_1\Entails !A$ و هم $\Gamma_2\Entails\lnot !A$. آنگاه
$\Gamma_1\cup\Gamma_2$ سازگار است.
\end{thm}

\begin{proof}
خلاف ادعا فرض کنید $\Gamma_1\cup\Gamma_2$ ناسازگار است. آنگاه
$\Gamma_1,\Gamma_2\Sequent\ $ یک !!a{proof} دارد. لم مائهارا را بر
$\maeh{\Gamma_1}{\Gamma_2}\Sequent\maeh{\ }{\ }$ اعمال کنید تا
!!a{sentence}~$!A$ای در !!{language}
$\Lang L(\Gamma_1)\cap\Lang L(\Gamma_2)$ به دست آید، به‌طوری‌که
$\Proves\Gamma_1\Sequent !A$ و
$\Proves !A,\Gamma_2\Sequent\quad$. از دومی
$\Proves\Gamma_2\Sequent\lnot !A$ به دست می‌آید، اما فرض کرده بودیم
چنین !!{sentence}‌ای وجود ندارد.
\end{proof}

در بالا به‌طور ضمنی فرض کردیم نظریه‌های $\Gamma$، $\Gamma_1$ و
$\Gamma_2$ متناهی‌اند. با فشردگی، نتایج برای نظریه‌های نامتناهی نیز
برقرارند.

\end{document}
